\chapter*{Abstract}
\label{chap:abstract}
\addcontentsline{toc}{chapter}{\nameref{chap:abstract}}

Multi-camera person tracking is essential for security and operational monitoring in large-scale environments such as university campuses and industrial facilities. Existing approaches rely entirely on appearance-based re-identification (Re-ID) for cross-camera association, which is computationally expensive and degrades when subjects wear similar clothing, as is common in factory settings. We propose \usevar{\srsTitle}, a real-time multi-camera person tracking system that addresses these limitations through a two-tier identity matching strategy. For cameras with overlapping fields of view, identities are resolved geometrically using a homography-based spatial matching module that projects detected foot positions onto a shared 2D floor plan; matching is performed via Euclidean distance minimization using the Hungarian algorithm. For non-overlapping blind spots, an appearance-based Re-ID module using a YOLOv26m detector, ByteTrack tracker, and OSNet-AIN backbone is triggered only when a tracked person enters a 5\% camera-edge zone, extracting 512-dimensional embeddings for cosine similarity matching. The system was evaluated on the MTMMC dataset across a university campus and an industrial factory, achieving an average Multiple Object Tracking Accuracy (MOTA) of 54.3\%, IDF1 of 61.8\%, and a per-frame processing latency of 49.4 ms across four concurrent camera streams. Future work will focus on automating homography calibration and validating the system on additional real-world multi-camera benchmarks.

\par
\textit{\textbf{Keywords}---Multi-Camera Tracking, Person Re-Identification, Computer Vision, Situational Awareness, Security Automation}